\documentclass[11pt,a4paper]{article}
\usepackage{fontspec,amsmath,amssymb,amsthm,microtype}
\setmainfont{lmroman10-regular.otf}[BoldFont=lmroman10-bold.otf,ItalicFont=lmroman10-italic.otf,BoldItalicFont=lmroman10-bolditalic.otf]
\usepackage{xcolor}
\usepackage[most]{tcolorbox}
\usepackage{mathrsfs}
\providecommand{\lt}{<}
\providecommand{\gt}{>}
\setmonofont{DejaVuSansMono.ttf}[BoldFont=DejaVuSansMono-Bold.ttf,ItalicFont=DejaVuSansMono-Oblique.ttf,BoldItalicFont=DejaVuSansMono-BoldOblique.ttf,Scale=0.86]
\newfontfamily\leanfallback{STIXTwoMath-Regular.otf}[Scale=0.86]
\definecolor{leanrule}{gray}{0.62}
\definecolor{leanback}{gray}{0.962}
\definecolor{leantitle}{gray}{0.22}
\definecolor{leanmuted}{gray}{0.38}
\newtcolorbox{leanbox}[1]{enhanced,breakable,colback=leanback,colframe=leanrule,
  boxrule=0.4pt,arc=0.6mm,left=2.4mm,right=2.4mm,top=1.8mm,bottom=1.8mm,
  before skip=0.8em,after skip=1em,fonttitle=\small\bfseries,coltitle=leantitle,
  colbacktitle=leanback,titlerule=0.3pt,title={#1}}
\newlength{\leancharwidth}
\newenvironment{leancode}{\par\smallskip\begingroup\ttfamily\footnotesize\raggedright
  \settowidth{\leancharwidth}{M}\setlength{\parskip}{0pt}\setlength{\parindent}{0pt}
  \frenchspacing}{\par\endgroup\smallskip}
\newcommand{\leanline}[2]{\hangindent=\dimexpr#1\leancharwidth+4\leancharwidth\relax
  \hangafter=1\noindent\hspace*{#1\leancharwidth}#2\par}
\newcommand{\leansource}[1]{\par{\footnotesize\color{leanmuted}#1}\par}
\newcommand{\leanlabel}[1]{\par\smallskip\noindent{\small\bfseries #1}\par}
\usepackage[margin=25mm]{geometry}
\usepackage[unicode,hidelinks]{hyperref}
\hypersetup{pdfauthor={LangevinLean Project}}
\setlength{\emergencystretch}{3em}
\setlength{\parskip}{0.4em}
\setlength{\parindent}{0pt}
\allowdisplaybreaks[2]
\raggedbottom
\title{\texorpdfstring{The four correlated Gaussian responses used by KLMC2}{The four correlated Gaussian responses used by KLMC2}}
\author{LangevinLean Project}
\date{}
\begin{document}\maketitle
Draft mathematical exposition; these notes may contain mistakes.


The second-order method uses four responses of the same Brownian interval: the two exponential-integrator kernels and the two Hessian correction kernels. Their correlations are essential. This section constructs those four responses simultaneously in every spatial coordinate, proves their exact Gram Gaussian law and iid time stream, and lifts the same random variables to arbitrary initial coupling extensions.



Four response kernels are evaluated on the same Brownian interval and packed into spatial vectors. Their full correlated block law and independence across intervals identify the second-order innovation stream.


\section{\texorpdfstring{Pack the exact scalar covariance into spatial vectors}{Pack the exact scalar covariance into spatial vectors}}


The Euclidean space uses an arbitrary finite spatial coordinate set. Packing transposes the array of scalar four-dimensional blocks into four spatial vectors. Every coordinate projection and Euclidean packing map is measurable. Pushing the product of the four-dimensional Gram Gaussians through this pack map defines the exact vector innovation law, which is a probability measure.



Each scalar block contains four correlated responses from one spatial Brownian coordinate. Packing exchanges the response and spatial indices without changing those values. The resulting four vectors are correlated exactly as required by the Hessian-corrected formula.

\begin{leanbox}{Lean: \texttt{H}, \texttt{pack}, \texttt{measurable\_\allowbreak{}pack} and 2 more}
\leansource{\texttt{H} \quad KineticLangevin/\allowbreak{}CanonicalKLMC2Noise.lean, lines 9–9}
\begin{leancode}
\leanline{0}{abbrev\ H\ :=\ EuclideanSpace\ ℝ\ ι}
\end{leancode}
\leansource{\texttt{pack} \quad KineticLangevin/\allowbreak{}CanonicalKLMC2Noise.lean, lines 12–13}
\begin{leancode}
\leanline{0}{def\ pack\ (x\ :\ ι\ →\ EuclideanSpace\ ℝ\ (Fin\ 4))\ :\ Fin\ 4\ →\ H\ ι\ :=}
\leanline{2}{fun\ j\ =>\ WithLp.toLp\ 2\ (fun\ i\ =>\ x\ i\ j)}
\end{leancode}
\leansource{\texttt{measurable\_\allowbreak{}pack} \quad KineticLangevin/\allowbreak{}CanonicalKLMC2Noise.lean, lines 15–19}
\begin{leancode}
\leanline{0}{theorem\ measurable\_pack\ :\ Measurable\ (pack\ ι)}
\end{leancode}
\leansource{\texttt{gaussianInnovationLaw} \quad KineticLangevin/\allowbreak{}CanonicalKLMC2Noise.lean, lines 22–23}
\begin{leancode}
\leanline{0}{def\ gaussianInnovationLaw\ (γ\ h\ :\ ℝ)\ :\ Measure\ (Fin\ 4\ →\ H\ ι)\ :=}
\leanline{2}{(Measure.pi\ (fun\ \_\ :\ ι\ =>\ klmc2Gaussian\ γ\ h)).map\ (pack\ ι)}
\end{leancode}
\leansource{\texttt{gaussianInnovationLaw\_\allowbreak{}isProbabilityMeasure} \quad KineticLangevin/\allowbreak{}CanonicalKLMC2Noise.lean, lines 25–27}
\begin{leancode}
\leanline{0}{instance\ gaussianInnovationLaw\_isProbabilityMeasure\ (γ\ h\ :\ ℝ)\ :}
\leanline{4}{IsProbabilityMeasure\ (gaussianInnovationLaw\ ι\ γ\ h)\ :=}
\leanline{2}{(Measure.isProbabilityMeasure\_map\_iff\ (measurable\_pack\ ι).aemeasurable).mpr\ inferInstance}
\end{leancode}
\leanlabel{How the Lean proof follows the argument}
\texttt{H} names the mathematical object defined above. Its defining expression appears in the adjacent statement; subsequent identities may unfold this definition directly.




\texttt{pack} names the mathematical object defined above. Its defining expression appears in the adjacent statement; subsequent identities may unfold this definition directly.




\texttt{measurable\_\allowbreak{}pack} gives the formal counterpart of this step. Named proof applications include \texttt{measurable\_\allowbreak{}pi\_\allowbreak{}iff.\allowbreak{}mpr}; the adjacent proof shows their hypotheses and order of use.




\texttt{gaussian\allowbreak{}Innovation\allowbreak{}Law} names the mathematical object defined above. Its defining expression appears in the adjacent statement; subsequent identities may unfold this definition directly.




\texttt{gaussian\allowbreak{}Innovation\allowbreak{}Law\_\allowbreak{}is\allowbreak{}Probability\allowbreak{}Measure} names the mathematical object defined above. Its defining expression appears in the adjacent statement; subsequent identities may unfold this definition directly.
\end{leanbox}

\section{\texorpdfstring{The Brownian innovation stream}{The Brownian innovation stream}}


Apply all four deterministic kernels to the same Brownian interval and pack all independent spatial copies. Brownian intervals have disjoint supports, so these full correlated vectors are independent across time. At each step their law is the packed product of the exact scalar Gram Gaussian laws. The independent-family product-law theorem therefore proves the complete defined stream has the iid Gaussian distribution.



The four responses on one interval remain dependent, while blocks from disjoint intervals are independent. This distinction must be preserved when specifying the stream law. The complete product law supplies fresh vector blocks at each numerical step.

\begin{leanbox}{Lean: \texttt{innovation}, \texttt{measurable\_\allowbreak{}innovation}, \texttt{innovation\_\allowbreak{}independent} and 2 more}
\leansource{\texttt{innovation} \quad KineticLangevin/\allowbreak{}CanonicalKLMC2Noise.lean, lines 30–31}
\begin{leancode}
\leanline{0}{def\ innovation\ (γ\ :\ ℝ)\ (h\ :\ ℝ≥0)\ (n\ :\ ℕ)\ (ω\ :\ ULA.CanonicalBrownian.VectorSample\ ι)\ :\ Fin\ 4\ →\ H\ ι\ :=}
\leanline{2}{pack\ ι\ (vectorGridBlock\ ι\ (klmc2Kernels\ γ)\ (continuous\_klmc2Kernels\ γ)\ h\ n\ ω)}
\end{leancode}
\leansource{\texttt{measurable\_\allowbreak{}innovation} \quad KineticLangevin/\allowbreak{}CanonicalKLMC2Noise.lean, lines 33–34}
\begin{leancode}
\leanline{0}{theorem\ measurable\_innovation\ (γ\ :\ ℝ)\ (h\ :\ ℝ≥0)\ (n\ :\ ℕ)\ :\ Measurable\ (innovation\ ι\ γ\ h\ n)}
\end{leancode}
\leansource{\texttt{innovation\_\allowbreak{}independent} \quad KineticLangevin/\allowbreak{}CanonicalKLMC2Noise.lean, lines 37–40}
\begin{leancode}
\leanline{0}{theorem\ innovation\_independent\ (γ\ :\ ℝ)\ (h\ :\ ℝ≥0)\ :}
\leanline{4}{iIndepFun\ (innovation\ ι\ γ\ h)\ (ULA.CanonicalBrownian.vectorProbability\ ι)}
\end{leancode}
\leansource{\texttt{innovation\_\allowbreak{}law} \quad KineticLangevin/\allowbreak{}CanonicalKLMC2Noise.lean, lines 43–54}
\begin{leancode}
\leanline{0}{theorem\ innovation\_law\ (γ\ :\ ℝ)\ (h\ :\ ℝ≥0)\ (n\ :\ ℕ)\ :}
\leanline{4}{HasLaw\ (innovation\ ι\ γ\ h\ n)\ (gaussianInnovationLaw\ ι\ γ\ h)\ (ULA.CanonicalBrownian.vectorProbability\ ι)}
\end{leancode}
\leansource{\texttt{innovation\_\allowbreak{}streamLaw} \quad KineticLangevin/\allowbreak{}CanonicalKLMC2Noise.lean, lines 57–62}
\begin{leancode}
\leanline{0}{theorem\ innovation\_streamLaw\ (γ\ :\ ℝ)\ (h\ :\ ℝ≥0)\ :}
\leanline{4}{HasLaw\ (fun\ ω\ n\ =>\ innovation\ ι\ γ\ h\ n\ ω)}
\leanline{6}{(Measure.infinitePi\ (fun\ \_\ :\ ℕ\ =>\ gaussianInnovationLaw\ ι\ γ\ h))}
\leanline{6}{(ULA.CanonicalBrownian.vectorProbability\ ι)}
\end{leancode}
\leanlabel{How the Lean proof follows the argument}
\texttt{innovation} names the mathematical object defined above. Its defining expression appears in the adjacent statement; subsequent identities may unfold this definition directly.




\texttt{measurable\_\allowbreak{}innovation} gives the formal counterpart of this step.




\texttt{innovation\_\allowbreak{}independent} gives the formal counterpart of this step.




\texttt{innovation\_\allowbreak{}law} gives the formal counterpart of this step. The main intermediate facts are \texttt{R}, \texttt{hb}, \texttt{hp} in that order. Named proof applications include \texttt{hp.\allowbreak{}comp}; the adjacent proof shows their hypotheses and order of use.




\texttt{innovation\_\allowbreak{}stream\allowbreak{}Law} gives the formal counterpart of this step.
\end{leanbox}


Earlier results used here: \cite[Section~3.8.16]{LangevinLean2026ULASC}; \cite[Section~3.8.17]{LangevinLean2026ULASC}.

\section{\texorpdfstring{The \(L^{2}\) variables with the original sampler noise}{The L² variables with the original sampler noise}}


The four \(L^{2}\) innovations are directly the vector Brownian convolutions, so no separate Gaussian realization is substituted. Their representative theorem identifies every coordinate with the corresponding defined packed response on the same source. Finite second moments follow from the \(L^{2}\) construction and this exact almost-everywhere identity.



The representative identification is pathwise almost everywhere and includes all four innovations. Hence cancellation involving Hessian-multiplied noise is valid on the coupling source. Merely identifying their separate marginal distributions would not establish that cancellation.

\begin{leanbox}{Lean: \texttt{innovationLp}, \texttt{innovationLp\_\allowbreak{}coe}, \texttt{innovation\_\allowbreak{}memLp}}
\leansource{\texttt{innovationLp} \quad KineticLangevin/\allowbreak{}CanonicalKLMC2Noise.lean, lines 65–67}
\begin{leancode}
\leanline{0}{def\ innovationLp\ (γ\ :\ ℝ)\ (h\ :\ ℝ≥0)\ (n\ :\ ℕ)\ (j\ :\ Fin\ 4)\ :}
\leanline{4}{Lp\ (H\ ι)\ 2\ (ULA.CanonicalBrownian.vectorProbability\ ι)\ :=}
\leanline{2}{vectorSegmentLp\ ι\ (klmc2Kernels\ γ\ j)\ (continuous\_klmc2Kernels\ γ\ j)\ ((n\ :\ ℝ≥0)*h)\ h}
\end{leancode}
\leansource{\texttt{innovationLp\_\allowbreak{}coe} \quad KineticLangevin/\allowbreak{}CanonicalKLMC2Noise.lean, lines 70–76}
\begin{leancode}
\leanline{0}{theorem\ innovationLp\_coe\ (γ\ :\ ℝ)\ (h\ :\ ℝ≥0)\ (n\ :\ ℕ)\ :}
\leanline{4}{(fun\ ω\ j\ =>\ innovationLp\ ι\ γ\ h\ n\ j\ ω)\ =ᵐ[ULA.CanonicalBrownian.vectorProbability\ ι]\ innovation\ ι\ γ\ h\ n}
\end{leancode}
\leansource{\texttt{innovation\_\allowbreak{}memLp} \quad KineticLangevin/\allowbreak{}CanonicalKLMC2Noise.lean, lines 79–84}
\begin{leancode}
\leanline{0}{theorem\ innovation\_memLp\ (γ\ :\ ℝ)\ (h\ :\ ℝ≥0)\ (n\ :\ ℕ)\ (j\ :\ Fin\ 4)\ :}
\leanline{4}{MemLp\ (fun\ ω\ =>\ innovation\ ι\ γ\ h\ n\ ω\ j)\ 2\ (ULA.CanonicalBrownian.vectorProbability\ ι)}
\end{leancode}
\leanlabel{How the Lean proof follows the argument}
\texttt{innovation\allowbreak{}Lp} names the mathematical object defined above. Its defining expression appears in the adjacent statement; subsequent identities may unfold this definition directly.




\texttt{innovation\allowbreak{}Lp\_\allowbreak{}coe} gives the formal counterpart of this step. The main intermediate facts are \texttt{hs} in that order. Named proof applications include \texttt{h\ensuremath{\omega}}; the adjacent proof shows their hypotheses and order of use.




\texttt{innovation\_\allowbreak{}mem\allowbreak{}Lp} gives the formal counterpart of this step.
\end{leanbox}


Earlier results used here: \cite[Section~3.8.17]{LangevinLean2026ULASC}.

\section{\texorpdfstring{The noise on every initial coupling extension}{The noise on every initial coupling extension}}


Pull the same Brownian convolutions onto the initial-position-coupling/\allowbreak{}Gaussian-velocity/\allowbreak{}Brownian product. The Brownian source projection preserves measure. The resulting representatives are exactly the original innovation values at that coordinate. A countable almost-everywhere intersection identifies the entire lifted stream with the original stream, and source projection preserves its exact iid law.



Adjoining the initial transport plan and shared velocity preserves the Brownian source law through its projection. Pulling back the responses therefore retains both their second moments and their full joint stream distribution. The local analysis uses this same lifted source.

\begin{leanbox}{Lean: \texttt{couplingProbability}, \texttt{extensionInnovation}, \texttt{extensionInnovation\_\allowbreak{}coe} and 1 more}
\leansource{\texttt{couplingProbability} \quad KineticLangevin/\allowbreak{}CanonicalKLMC2Noise.lean, lines 87–87}
\begin{leancode}
\leanline{0}{local\ instance\ couplingProbability\ :\ IsProbabilityMeasure\ c.joint\ :=\ c.isProbabilityMeasure}
\end{leancode}
\leansource{\texttt{extensionInnovation} \quad KineticLangevin/\allowbreak{}CanonicalKLMC2Noise.lean, lines 90–92}
\begin{leancode}
\leanline{0}{def\ extensionInnovation\ (γ\ :\ ℝ)\ (h\ :\ ℝ≥0)\ (n\ :\ ℕ)\ (j\ :\ Fin\ 4)\ :}
\leanline{4}{Lp\ (H\ ι)\ 2\ (GenericCanonicalCoupling.probability\ c)\ :=}
\leanline{2}{brownianLift\ c\ (innovationLp\ ι\ γ\ h\ n\ j)}
\end{leancode}
\leansource{\texttt{extensionInnovation\_\allowbreak{}coe} \quad KineticLangevin/\allowbreak{}CanonicalKLMC2Noise.lean, lines 95–106}
\begin{leancode}
\leanline{0}{theorem\ extensionInnovation\_coe\ (γ\ :\ ℝ)\ (h\ :\ ℝ≥0)\ (n\ :\ ℕ)\ :}
\leanline{4}{(fun\ p\ j\ =>\ extensionInnovation\ c\ γ\ h\ n\ j\ p)\ =ᵐ[GenericCanonicalCoupling.probability\ c]}
\leanline{6}{fun\ p\ :\ GenericCanonicalCoupling.Sample\ ι\ =>\ innovation\ ι\ γ\ h\ n\ p.2.2}
\end{leancode}
\leansource{\texttt{extensionInnovation\_\allowbreak{}streamLaw} \quad KineticLangevin/\allowbreak{}CanonicalKLMC2Noise.lean, lines 109–121}
\begin{leancode}
\leanline{0}{theorem\ extensionInnovation\_streamLaw\ (γ\ :\ ℝ)\ (h\ :\ ℝ≥0)\ :}
\leanline{4}{HasLaw\ (fun\ p\ n\ j\ =>\ extensionInnovation\ c\ γ\ h\ n\ j\ p)}
\leanline{6}{(Measure.infinitePi\ (fun\ \_\ :\ ℕ\ =>\ gaussianInnovationLaw\ ι\ γ\ h))}
\leanline{6}{(GenericCanonicalCoupling.probability\ c)}
\end{leancode}
\leanlabel{How the Lean proof follows the argument}
\texttt{coupling\allowbreak{}Probability} names the mathematical object defined above. Its defining expression appears in the adjacent statement; subsequent identities may unfold this definition directly.




\texttt{extension\allowbreak{}Innovation} names the mathematical object defined above. Its defining expression appears in the adjacent statement; subsequent identities may unfold this definition directly.




\texttt{extension\allowbreak{}Innovation\_\allowbreak{}coe} gives the formal counterpart of this step. The main intermediate facts are \texttt{hp}, \texttt{he}, \texttt{hs} in that order.




\texttt{extension\allowbreak{}Innovation\_\allowbreak{}stream\allowbreak{}Law} gives the formal counterpart of this step. The main intermediate facts are \texttt{hp}, \texttt{hl}, \texttt{hs} in that order. Named proof applications include \texttt{hl.\allowbreak{}congr}, \texttt{funext}; the adjacent proof shows their hypotheses and order of use.
\end{leanbox}


Earlier results used here: \cite[Section~3.8.17]{LangevinLean2026ULASC}.

\begin{thebibliography}{ULA-LSI}
\bibitem[ULA-SC]{LangevinLean2026ULASC}
LangevinLean Project. \emph{Unadjusted Langevin under Strong Convexity}. LL-01, 2026. Draft report, content revision d50c0f64815b. \url{https://odakyildiz.github.io/LangevinLean/docs/pdf/UnadjustedLangevin/FullProof.pdf}.
\end{thebibliography}

\end{document}
